% !TeX program = xelatex
\documentclass[UTF8,11pt,fontset=windows]{ctexart}
\usepackage[a4paper,left=22mm,right=22mm,top=20mm,bottom=20mm]{geometry}
\usepackage{fontspec}
\setmainfont{Times New Roman}
\setsansfont{Arial}
\usepackage{array,booktabs,tabularx}
\usepackage{xcolor,enumitem,amsmath,fancyhdr,hyperref}
\hypersetup{colorlinks=true,linkcolor=black,citecolor=blue!45!black,urlcolor=blue!45!black,pdftitle={浙工新村模式与经济机制研究},pdfauthor={朱航成}}
\renewcommand{\arraystretch}{1.23}
\setlength{\tabcolsep}{4pt}
\setlength{\parskip}{3pt}
\linespread{1.15}
\setlist[itemize]{leftmargin=1.8em,itemsep=2pt,topsep=3pt}
\ctexset{section={format=\large\bfseries,beforeskip=10pt,afterskip=5pt},subsection={format=\normalsize\bfseries,beforeskip=6pt,afterskip=4pt}}
\pagestyle{fancy}\fancyhf{}\fancyfoot[C]{\small\thepage}
\renewcommand{\headrulewidth}{0pt}\setlength{\headheight}{14pt}
\raggedbottom\clubpenalty=10000\widowpenalty=10000\emergencystretch=2em
\newcolumntype{L}[1]{>{\raggedright\arraybackslash}p{#1}}
\newcolumntype{Y}{>{\raggedright\arraybackslash}X}
\newcommand{\src}[1]{\cite{#1}}
\title{\bfseries\Large 浙工新村：模式组合与经济机制\\\normalsize 案例拆分、成本收益与研判建议}
\author{\normalsize 朱航成}
\date{\small 2026年10月12日}
\begin{document}
\maketitle\vspace{-6mm}

本研究将杭州拱墅区浙工新村拆解为具体措施、实施条件、资金流和市场效果，形成候选小区的对照材料。案例已完成重建，当前研究方向为运行评估与资金复盘；对其他区域的建议按模式准入、经济比较和实施安排逐层形成。

\section{模式先判，再选择措施组合}
浙工新村采用业主自主更新和危旧房原址重建，并组合户型与面积再配置、地下停车、公共服务和治理运营。2024年方案记载，新建7幢11层住宅、548套，住宅套数保持548套；总建筑面积约8.1万平方米，其中地上约5.7万、地下约2.4万平方米。\src{S04}

\noindent\begin{tabularx}{\textwidth}{@{}L{30mm}YY@{}}
\toprule
\textbf{措施单元} & \textbf{案例怎么做} & \textbf{应提取的关键条件}\\
\midrule
危旧房原址重建 & 4幢房屋曾被鉴定为C级危房；纳入试点范围的原楼整体重建。\src{S03} & 检测结论、维修加固效果、原址意愿、权籍、规划及实施范围。\\
面积与户型优化 & 保持原套内面积，自选有偿扩面上限20平方米/户。\src{S04} & 原建筑/套内面积、户型、付费扩面、面积分配与居民预算。\\
地下停车 & 新建地下车库，人车分流，车位供居民自选购买。\src{S04,S05} & 地下空间条件、停车需求、销售权属、专项成本和实际回款。\\
服务与环境 & 配置一老一小等公共空间，无障碍设施和北侧沿街商铺。\src{S03,S06} & 服务人群、面积用途、产权、运营方、收费和持续维护。\\
协商与管理 & 居民参与协商、审计及质量沟通，交付后进入物业管理。\src{S05,S08} & 决议、施工监督、实收与结算、维护责任、居民反馈。\\
\bottomrule
\end{tabularx}

\subsection{迁移时比较多种更新组合}
候选区域可同时建立五组方案：保留综合整治、加固与功能更新、局部拆建加保留整治、整体自主拆建、居住更新加经营服务空间。物理更新、治理组织和资金包分别配置；逐幢检测与权籍资料决定能选什么，居民需求和经济条件决定优先做什么。

\textbf{知识沉淀字段：}对象、动作、条件、工程量、费用口径、产权及受益方、实施阶段、结果与维护。每项字段连接原文位置、形成时点和版本；计划、施工、交付和运行分别保存。

\newpage
\section{两版资金参数与各主体的账}
\noindent\begin{tabularx}{\textwidth}{@{}L{36mm}L{29mm}L{31mm}Y@{}}
\toprule
\textbf{报道版本} & \textbf{整体费用} & \textbf{居民自筹} & \textbf{资金组织}\\
\midrule
2024年1月19日 & 约5.3亿元 & 约4.7亿元 & 其余由旧改、加梯、未来社区等专项支持衔接。\src{S03}\\
2026年8月19日回顾 & 约5.6亿元 & 约4.9亿元 & 国企在建设期垫资，回迁结算；不足部分由专项资金补齐。\src{S08}\\
\bottomrule
\end{tabularx}

以上分别保存为历史预算和回顾报道版本；分项终审决算、资金实收与借款合同列入采集清单。按各版总额计算，居民出资占比分别约88.7\%和87.5\%。模型还要读取每笔款项的支付方、收款方、用途和到账日期。

\subsection{造血机制：付款、收款与受益分别记录}
\small
\noindent\begin{tabularx}{\textwidth}{@{}L{24mm}L{27mm}YY@{}}
\toprule
\textbf{渠道} & \textbf{付款方} & \textbf{收款/权利主体} & \textbf{对应经济效果}\\
\midrule
原面积与扩面出资 & 原产权家庭 & 建设结算主体，按合同核对 & 筹集建设资金；家庭获得新住房和自选面积。\\
专项支持 & 财政/专项渠道 & 项目主体或居民，按分项核对 & 支持公共改善、过渡及相关建设安排。\\
车位销售 & 购位业主 & 车位产权/销售主体 & 资产回收，居民获得车位；净额扣对应建设及费用。\\
商铺租赁与服务 & 承租经营者/服务付费人 & 产权人或合法运营方 & 潜在长期现金流，扣空置、税费、维护和运营。\\
自住资产变化 & 持有时无付款方 & 原产权家庭持有 & 更新后的居住改善及住房资产变化。\\
原产权二手出售 & 二手买家 & 出售房屋的原产权人 & 家庭售房回款，另核税费、剩余贷款和此前投入。\\
获批新增外售住房 & 新增住房买家 & 有可售权属的主体 & 按批准面积去向与销售条件设候选收入项。\\
\bottomrule
\end{tabularx}
\normalsize

浙工新村的建设平衡首先依靠居民出资与公共支持；家庭改善和房产价值变化体现居民收益，二手成交体现产权人回款。企业或运营方的收入由承包、代建、融资、合法资产销售及运营合同确定。每个方案分别展示建设资金、运营收支和居民经济效果，便于判断各主体能否持续参与。

\newpage
\section{市场需求、家庭负担与收益测算}
\small
\noindent\begin{tabularx}{\textwidth}{@{}L{33mm}L{39mm}Y@{}}
\toprule
\textbf{参数/样本} & \textbf{数值与时点} & \textbf{模型用途及资料范围}\\
\midrule
原面积出资 & 1350元/平方米，2025年2月 & 原住房计价面积，不足53平方米按53平方米计。\src{S05}\\
自选扩面出资 & 34520元/平方米，2025年2月 & 专业评估后的扩面价格；每户最多自选20平方米。\src{S05}\\
地下车位 & 22.1万元/个，2025年2月 & 自愿购买参数；销售量、权属及专项成本另采。\src{S05}\\
首套二手网签 & 41114元/平方米，2025年5月上旬 & 杭州贝壳研究院的一套样本；5-6月成交15套。\src{S07}\\
更新后市场观察 & 超过4万元/平方米，2026年8月 & 采访概括；交付后成交100多套，补逐笔匹配样本。\src{S08}\\
原业主与改善买家 & 自住养老及周边改善需求 & 分别调查基本改造、扩面、车位支付；二手买家比较同总价替代住房。\src{S05}\\
\bottomrule
\end{tabularx}
\normalsize

\subsection{一个家庭算例与三个计算账户}
以计价原面积75平方米、自选扩面20平方米、购买1个车位为假设，历史报价对应：
\[
75\times1350+20\times34520+221000=101.265\text{万元}
\]
不购车位为79.165万元，补助、装修、租住过渡与融资分别输入。该算例用于比较家庭方案，并可改变扩面和车位参数观察支付压力。

\small
\noindent\begin{tabularx}{\textwidth}{@{}L{28mm}Y@{}}
\toprule
\textbf{账户} & \textbf{计算内容}\\
\midrule
建设及运营 & 居民实收＋政府支持实收＋可售资产实收＋经营净现金，扣工程、过渡、咨询迁改、税费、融资与承担的维护。\\
居民净资产变化 & 更新后同类可实现住房价值－更新前同类住房价值－家庭净新增投入－交易与融资成本。\\
居民出售回款 & 二手出售实收－税费－剩余贷款，另附此前追加投入记录。\\
\bottomrule
\end{tabularx}
\normalsize

\subsection{缺项与采集触发}
真实租约、商铺面积与经营权、车位实收及专项成本、垫资利率、分户支付能力、年度维护预算均保留\texttt{null / missing}。渠道分别为租约与产权档案、销售和成本台账、银行合同、分层问卷与物业预算。拟用商铺与停车净收入承担维护而上述字段缺失时，输出“先补充关键资料再判断”，生成对应采集任务，再更新收益和模式比较。

\newpage
\section{八项研究决策及其依据}
\small
\noindent\begin{tabularx}{\textwidth}{@{}L{10mm}Y L{53mm}@{}}
\toprule
\textbf{编号} & \textbf{建议与执行内容} & \textbf{依据及对应规则}\\
\midrule
D1 & 将浙工新村转入运行评估与资金复盘，汇总交房、产权、实收和运维资料。 & 已重建交付和竣工投用；阶段匹配。\src{S06,S10}\\
D2 & 为新区域比较保留、加固、局部与整体拆建，以及服务经营空间组合。 & 逐幢安全、权籍与功能问题；模式准入和规划条件。\src{S03,S09}\\
D3 & 分开原面积基础包、扩面包和车位包，按家庭支付能力验证。 & 历史收费及交付付款；支付负担和面积分配。\src{S05,S06}\\
D4 & 建设、运营、居民三账分列，关联权属、收款与成本。 & 居民筹资、国企垫资、车位及商铺；资金平衡。\src{S03,S05,S06,S08}\\
D5 & 用当地可比成交构建保守、基准、上行情景，分别报资产变化和经营现金。 & 网签及成交资料；同品质可比、需求与去化。\src{S07,S08}\\
D6 & 对缺少租约、权属和专项成本的经营方案生成补充调查任务。 & 沿街商铺及车位的收入字段；收益条件和数据覆盖。\src{S05,S06}\\
D7 & 按月模拟国企垫资、居民结算和补助到账，计算资金峰值。 & 建设期垫资与交付付清；支付时序和流动性。\src{S06,S08}\\
D8 & 以安全、功能、居民负担、建设收支、维护持续共同比较方案。 & 规划、居民决议、维护制度；硬条件与目标比较。\src{S09}\\
\bottomrule
\end{tabularx}
\normalsize

\subsection{对应政策规则与专家复核问题}
以下来自2024年浙江自主更新试点文件，作为规则候选附原文位置；新候选区域同时核对所在市县的现行适用口径。\src{S09}
\begin{itemize}
\item \textbf{R1 组织与产权：}更新方案包括组织、建设代理、筹资、产权处置和设计（正文第4-5页）。
\item \textbf{R2 面积与成本：}新增面积按地方政策确定去向，详细规划依法调整；保障房或服务设施移交可衔接成本（第8页）。
\item \textbf{R3 安全与规划：}保障公共安全，对间距、日照等进行技术处理，以现状条件为底线（第8页）。
\item \textbf{R4 长期维护：}专业物业管理和维修资金归集、使用、续筹形成持续运行安排（第9页）。
\end{itemize}
老师可进一步确定：检测结果怎样改变模式、支付能力怎样分组、公共与自选措施如何分摊、市场样本和经营权怎样进入收益计算、何种维护安排足以支撑长期运行。各阈值经复核后再进入版本化规则。

\newpage
\section{来源与字段追溯}
公开资料保留“来源日期、数据发生时点、空间与费用范围、原文定位”。资金两版本分列；13幢与14幢原楼报道分别连接实施范围清单，核准原面积、容积率与新增面积去向通过权籍及规划档案补齐。建议、计算假设和缺项与历史事实分别存储。

\vspace{2mm}
\small
\begin{thebibliography}{S10}
\bibitem[S01]{S01} 拱墅区政府·朝晖街道，2024-02-21。\href{https://www.gongshu.gov.cn/art/2024/2/21/art\_1229782350\_59079232.html}{打造大运河畔富有人间烟火气的幸福家园 开创街道高质量发展新局面}。定位：解锁亚运“时代密码”段；正文L31-L33。

\bibitem[S02]{S02} 拱墅区政府，2023-06-26。\href{https://www.gongshu.gov.cn/art/2023/6/26/art\_1230544\_59075321.html}{中共杭州市拱墅区委关于十五届省委第一轮巡视整改进展情况的通报}。定位：九、（二）全力推进浙工新村危房改造。

\bibitem[S03]{S03} 杭州政协网转载杭州日报，2024-01-19。\href{https://www.hzzx.gov.cn/cshz/content/2024-01/19/content\_8677402.htm}{全省率先试点城市危旧房有机更新 拱墅区如何做到“第一个吃螃蟹”？}。定位：正文L37-L57：工程规模、历史预算和住户算账。

\bibitem[S04]{S04} 杭州网转载潮新闻，2024-11-11。\href{https://hznews.hangzhou.com.cn/chengshi/content/2024-11/11/content\_8811527\_0.htm}{选房啦！杭州首个自主更新小区 传来好消息}。定位：正文L64-L66：重建方案与总建筑面积。

\bibitem[S05]{S05} 杭州网转载都市快报，2025-02-11。\href{https://hwyst.hangzhou.com.cn/mtld/content/2025-02/11/content\_8857599.htm}{杭州首个“原拆原建”自主更新的小区——浙工新村下月交付}。定位：原面积结算、扩面价格、车位及目标客群段（压缩正文L7）。

\bibitem[S06]{S06} 杭州网·每日商报，2025-02-27。\href{https://hznews.hangzhou.com.cn/chengshi/content/2025-02/27/content\_8868888.htm}{杭州首个“原拆原建”自住更新小区交付}。定位：正文L21-L32：距离、交付、沿街商铺与付款。

\bibitem[S07]{S07} 杭州网报道杭州贝壳研究院数据，2025-07-31。\href{https://hznews.hangzhou.com.cn/jingji/content/2025-07/31/content\_9050608\_2.htm}{杭州市中心两大焦点安置房逆势热销}。定位：正文L19-L21：首套网签及成交套数。

\bibitem[S08]{S08} 解放日报·上观新闻，2026-08-19。\href{https://www.jfdaily.com/wx/detail.do?id=1162310}{老旧小区原拆原建签约率100\%，居民自筹资金自主更新，其他地方能学吗？}。定位：正文L22-L24、L62、L95-L109、L149；报道日期见L7。

\bibitem[S09]{S09} 浙江省住房和城乡建设厅、发展改革委、自然资源厅，2024-04-01。\href{https://zjjcmspublic.oss-cn-hangzhou-zwynet-d01-a.internet.cloud.zj.gov.cn/jcms\_files/jcms1/web1838/site/attach/0/a66346da98264bbdbadf0310ebc148c6.pdf}{关于稳步推进城镇老旧小区自主更新试点工作的指导意见（试行）}。定位：PDF正文第4-9页：组织实施、政策保障、标准、管理。

\bibitem[S10]{S10} 今日拱墅，2025-02-28。\href{https://gsb.qybk.cn/resfile/2025-02-28/02/jrgs-20250228-002.pdf}{拱墅12个重大项目集中开工15个项目竣工}。定位：PDF第2版导语：浙工新村试点项目竣工投用。

\end{thebibliography}
\normalsize
\subsection{结构化案例包}
配套案例包保存36项事实、5张措施卡、11项缺项、8条研究建议和4组规则候选。状态字段：\texttt{factual} 为已核对政府文件或官方陈述；\texttt{reference} 为公开报道、采访和市场样本；\texttt{assumption} 为研究建议或计算情景；\texttt{missing} 为需要采集的字段。

每项资料新增后建立版本，重算其关联措施、资金账户和模式比较；来源和规则编号随建议输出，支持老师沿“建议、条件、原文、计算”逐项检查。
\end{document}
